%% 2-9 分数で書けない数がある？（第11回で書き足し。ページ数は自由・偶数）
%% この章で定義した部品（頭文字 bkiij）：
%%   \bkiijstrip{(左はし)}{左の文字}{右の文字}{色}   計算の答えを見せる横長の札
%%   \bkiijwheel{(中心)}{半径}                        自転車の車輪（スポークつき）
\makeatletter
\newcommand\bkiijstrip[4]{\begin{scope}[shift={#1}]
  \node[anchor=west, draw=#4!80!black, fill=#4!12, line width=0.7pt, rounded corners=1.5mm, minimum width=19mm, minimum height=7mm, font=\small, inner sep=1mm] (L) at (0,0) {#2};
  \node[anchor=west, font=\small, inner sep=0pt] at ([xshift=3mm]L.east) {#3};\end{scope}}
\newcommand\bkiijwheel[2]{\begin{scope}[shift={#1}]
  \draw[bkInk!85, line width=2.2pt] (0,0) circle (#2);
  \draw[bkGray!70, line width=0.8pt] (0,0) circle (#2-1.4);
  \foreach \a in {0,30,...,330} \draw[bkGray!80, line width=0.35pt] (0,0) -- (\a:#2-1.4);
  \fill[bkGray] (0,0) circle (1.2);\end{scope}}
% 小さな見出し  \bkiijhead{見出し}
\newcommand\bkiijhead[1]{\par\addvspace{5pt}\noindent{\sffamily\bfseries\color{bkTheme!85!black}▶\ #1}\par\nopagebreak\vspace{1pt}}
\makeatother

\bkchapter{2}{9}{\R{分数}{ぶん|すう}で\R{書}{か}けない\R{数}{かず}がある？}{\bklevel{2}}{\R{分数}{ぶん|すう}で\R{書}{か}けない\R{数}{かず}がある？}
\bkmichishirube{1-15,1-19,2-4,2-8}{2-19,4-2,4-4,5-14}

\begin{goal}
\begin{itemize}[label={\color{bkTheme}●}, leftmargin=1.6\zw, itemsep=0pt, topsep=1pt]
\item \bklevel{2}\enspace \R{分数}{ぶん|すう}を\R{小数}{しょう|すう}にすると、\R{終}{お}わるかくり\R{返}{かえ}すかのどちらかになるわけがわかる。
\item \bklevel{2}\enspace $\sqrt{2}$ や $\pi$ のように、\R{分数}{ぶん|すう}で\R{書}{か}けない\R{数}{かず}（\R{無理数}{む|り|すう}）があることがわかる。
\item \bklevel{2}\enspace \R{円周}{えん|しゅう}＝\R{直径}{ちょっ|けい}×$\pi$ を\R{使}{つか}って、\R{円周}{えん|しゅう}や\R{直径}{ちょっ|けい}を\R{求}{もと}められる。
\end{itemize}
\end{goal}

\begin{hatena}
\R{電卓}{でん|たく}で\R{計算}{けい|さん}してみましょう。$1\div4$ は0.25で\R{終}{お}わります。$1\div7$ は\R{終}{お}わりませんが、142857がくり\R{返}{かえ}されます。では、2-8の $\sqrt{2}$ は？　コンピュータで\R{何十}{なん|じゅう}けたも\R{計算}{けい|さん}しても、\R{終}{お}わりもしないし、くり\R{返}{かえ}しも\R{見}{み}つかりません。$\sqrt{2}$ は、\R{分数}{ぶん|すう}で\R{書}{か}けるのでしょうか？
\begin{bkfigure}
\begin{tikzpicture}[x=1mm, y=1mm]
  \bkiijstrip{(0,10)}{$1\div4$}{$0.25$\quad（\R{終}{お}わる）}{bkGreen}
  \bkiijstrip{(0,0)}{$1\div7$}{$0.\textcolor{bkOrange!90!black}{142857}\,142857\,142857\cdots$\quad（くり\R{返}{かえ}す）}{bkOrange}
  \bkiijstrip{(0,-10)}{$\sqrt{2}$}{$1.41421356237309504880168872420969807856\cdots$}{bkBlue}
\end{tikzpicture}
\end{bkfigure}
\end{hatena}

\R{分数}{ぶん|すう}は、わり\R{算}{ざん}の\R{答}{こた}えでした（2-4）。この\R{章}{しょう}では、\R{分数}{ぶん|すう}を\R{小数}{しょう|すう}にしたときのきまりを\R{見}{み}つけ、そのきまりにあてはまらない\R{数}{かず}をさがします。

\clearpage
%% ---------------- やってみよう・思い出そう ----------------
\begin{yattemiyou}[\R{円周}{えん|しゅう}÷\R{直径}{ちょっ|けい}をはかろう]
\noindent\begin{minipage}[c]{0.64\linewidth}
\begin{enumerate}[label={\bkcircled[bkGreen]{\arabic*}}, leftmargin=2.4\zw, itemsep=1pt, topsep=0pt]
\item かんやコップなど、\R{丸}{まる}い\R{物}{もの}を\R{用意}{よう|い}する。
\item \R{紙}{かみ}テープを1\R{周}{しゅう}まきつけ、\R{重}{かさ}なった\R{所}{ところ}に\R{印}{しるし}をつける。テープの\R{長}{なが}さが\textbf{\R{円周}{えん|しゅう}}。
\item \R{直径}{ちょっ|けい}をものさしではかる。
\item \R{電卓}{でん|たく}で、\R{円周}{えん|しゅう}÷\R{直径}{ちょっ|けい}を\R{計算}{けい|さん}しよう。
\end{enumerate}
\end{minipage}\hfill
\begin{minipage}[c]{0.33\linewidth}\centering
\begin{tikzpicture}[x=1mm, y=1mm]
  \fill[bkGray!20] (-10,0) rectangle (10,24);
  \fill[bkGray!30] (0,0) ellipse (10 and 3);
  \filldraw[fill=bkGray!12, draw=bkInk!60, line width=0.6pt] (0,24) ellipse (10 and 3);
  \draw[bkInk!60, line width=0.6pt] (-10,0) -- (-10,24) (10,0) -- (10,24);
  \draw[bkInk!60, line width=0.6pt] (-10,0) arc (180:360:10 and 3);
  \fill[bkYellow!70] (-10,10) arc (180:360:10 and 3) -- (10,15) arc (360:180:10 and 3) -- cycle;
  \draw[bkRed, line width=1pt] (3,7.2) -- (3,12.2);
  \draw[bkGray, line width=0.5pt, <->] (-10,30) -- (10,30);
  \node[anchor=south, font=\sffamily\footnotesize] at (0,31) {\R{直径}{ちょっ|けい}};
  \node[anchor=west, font=\sffamily\footnotesize, align=left] at (12,10) {\R{紙}{かみ}\\テープ};
\end{tikzpicture}
\end{minipage}
\par\vspace{4pt}
\begin{bkcenter}
\begin{bktable}{colspec={Q[l,m,wd=27mm] *{3}{Q[c,m,wd=17mm]} Q[c,m,wd=24mm]}, rowsep=2pt, colsep=4pt, row{1}={bg=bkGreen!14}, column{1}={font=\sffamily\small}}
\R{物}{もの} & かん & コップ & \R{皿}{さら} & テープのしん \\
\R{直径}{ちょっ|けい}（cm） & 6.6 & 7.5 & 20.0 & 7.6 \\
\R{円周}{えん|しゅう}（cm） & 20.8 & 23.5 & 62.8 & 23.9 \\
\R{円周}{えん|しゅう}÷\R{直径}{ちょっ|けい} & 3.15 & 3.13 & 3.14 & 3.14 \\
\end{bktable}
\bktabcaption{はかった\R{例}{れい}。\R{大}{おお}きさがちがっても、\R{円周}{えん|しゅう}÷\R{直径}{ちょっ|けい}は、いつも\R{同}{おな}じくらいの\R{数}{かず}になる。}
\end{bkcenter}
\end{yattemiyou}

\begin{omoidasou}
\bkrecall{2-4}{\R{分数}{ぶん|すう}はわり\R{算}{ざん}の\R{答}{こた}え。$2\div3=\dfrac{2}{3}$。\R{約分}{やく|ぶん}すると $\dfrac{25}{100}=\dfrac{1}{4}$}
\bkrecall{1-15}{わり\R{算}{ざん}のあまりは、わる\R{数}{かず}より\R{小}{ちい}さい。}
\bkrecall{2-8}{$\sqrt{2}$ は、2\R{乗}{じょう}すると2になる\R{正}{せい}の\R{数}{かず}。$\sqrt{2}=1.41421\cdots$}
\bkrecall{1-19}{\R{円}{えん}の\R{直径}{ちょっ|けい}は、\R{半径}{はん|けい}の2つ\R{分}{ぶん}。}
\end{omoidasou}

\clearpage
%% ---------------- 数学の道具 ----------------
\bksection{doguu}{\R{数学}{すう|がく}の\R{道具}{どう|ぐ}}

\bksubsection{1}{\R{分数}{ぶん|すう}を\R{小数}{しょう|すう}にすると、\R{終}{お}わるかくり\R{返}{かえ}す}

\leavevmode\label{bk:2-9:kurikaesu}\R{分数}{ぶん|すう}を\R{小数}{しょう|すう}にするには、\R{分子}{ぶん|し}÷\R{分母}{ぶん|ぼ}を\R{計算}{けい|さん}します（2-4）。まず、\R{終}{お}わる\R{例}{れい}の $\dfrac{1}{4}=1\div4$ を、1けたずつ\R{見}{み}てみましょう。

\bkiijhead{わり\R{算}{ざん}を1けたずつ\R{進}{すす}める}
1は4より\R{小}{ちい}さいので、\R{一}{いち}の\R{位}{くらい}は0です。1を0.1が10こ\R{分}{ぶん}と\R{考}{かんが}えて（2-1）、10を4でわります。$10\div4=2$ あまり2 なので、\R{小数第}{しょう|すう|だい}1\R{位}{い}は2。あまりの2（0.1が2こ）を、こんどは0.01が20こ\R{分}{ぶん}と\R{考}{かんが}えて、$20\div4=5$ あまり0。\R{小数第}{しょう|すう|だい}2\R{位}{い}は5で、あまりが0になったので\R{終}{お}わりです。$\dfrac{1}{4}=0.25$。

つまり、「\textbf{あまりに0をつけて10\R{倍}{ばい}し、また\R{分母}{ぶん|ぼ}でわる}」をくり\R{返}{かえ}すと、\R{小数}{しょう|すう}のけたが1つずつ\R{出}{で}てきます。\R{同}{おな}じように $\dfrac{1}{7}$ を、あまりに\R{注目}{ちゅう|もく}して\R{進}{すす}めてみましょう。

\begin{bkcenter}
\begin{bktable}{colspec={Q[l,m,wd=36mm] *{7}{Q[c,m,wd=8mm]}}, rowsep=1pt, colsep=2pt, cells={font=\small}, column{1}={font=\sffamily\small}, row{1}={font=\small}, cell{2}{2-7}={fg=bkOrange!90!black, font=\small\bfseries}}
わられる\R{数}{かず} & 10 & 30 & 20 & 60 & 40 & 50 & 10 \\
\R{商}{しょう}（\R{答}{こた}えのけた） & 1 & 4 & 2 & 8 & 5 & 7 & 1 \\
あまり & 3 & 2 & 6 & 4 & 5 & 1 & 3 \\
\end{bktable}
\end{bkcenter}
\noindent\begin{minipage}[c]{0.56\linewidth}
6\R{回}{かい}\R{目}{め}のあまりが1になり、\R{最初}{さい|しょ}と\R{同}{おな}じ\mbox{「10÷7」}にもどりました。\R{同}{おな}じわり\R{算}{ざん}からは、\R{同}{おな}じ\R{商}{しょう}と\R{同}{おな}じあまりが\R{出}{で}ます。だから、ここから\R{先}{さき}は\R{同}{おな}じ\R{計算}{けい|さん}のくり\R{返}{かえ}しで、
\[ \frac{1}{7}=0.142857\,142857\cdots \]
となります。\R{右}{みぎ}の\R{図}{ず}は、あまりがぐるぐる\R{回}{まわ}るようすです。
\end{minipage}\hfill
\begin{minipage}[c]{0.41\linewidth}\centering
\begin{tikzpicture}[x=1mm, y=1mm]
  \foreach \r/\q/\a [count=\i from 0] in {1/1/90, 3/4/30, 2/2/-30, 6/8/-90, 4/5/-150, 5/7/150} {
    \node[circle, draw=bkOrange!80!black, fill=bkOrange!12, line width=0.7pt, minimum size=7mm, inner sep=0pt, font=\small] (r\i) at (\a:17) {\r};}
  \foreach \i/\q [evaluate=\i as \j using {int(mod(\i+1,6))}] in {0/1,1/4,2/2,3/8,4/5,5/7} {
    \draw[bk link, bkOrange!80!black] (r\i) to[bend left=18] node[midway, font=\sffamily\footnotesize\bfseries, text=bkBlue, fill=white, inner sep=0.6pt] {\q} (r\j);}
  \node[font=\sffamily\footnotesize, align=center, inner sep=0pt] at (0,0) {\R{丸}{まる}：あまり\\\R{青}{あお}：\R{商}{しょう}};
\end{tikzpicture}
\end{minipage}

\bkiijhead{なぜ、かならずくり\R{返}{かえ}すの？}
7でわったあまりは、0〜6の7\R{通}{とお}りしかありません（1-15：あまりはわる\R{数}{かず}より\R{小}{ちい}さい）。
\begin{enumerate}[label={\bkcircled[bkGreen]{\arabic*}}, leftmargin=2.4\zw, itemsep=1pt, topsep=2pt]
\item あまりが0になれば、わり\R{算}{ざん}は\R{終}{お}わります（$\dfrac{1}{4}=0.25$ のように\R{終}{お}わる\R{小数}{しょう|すう}）。
\item 0にならないとき、あまりは1〜6の6\R{通}{とお}りのどれかです。わり\R{算}{ざん}を7\R{回}{かい}つづけると、あまりが7こ\R{出}{で}ます。6\R{通}{とお}りしかないのに7こ\R{出}{で}るので、どれか2つはかならず\R{同}{おな}じです（6つの\R{部屋}{へ|や}に7\R{人}{にん}が\R{入}{はい}ると、どこかの\R{部屋}{へ|や}は2\R{人}{にん}になるのと\R{同}{おな}じ）。
\item \R{同}{おな}じあまりが\R{出}{で}たら、そこから\R{先}{さき}は\R{同}{おな}じ\R{計算}{けい|さん}のくり\R{返}{かえ}しです。
\end{enumerate}
\R{分母}{ぶん|ぼ}が7でなくても、\R{考}{かんが}え\R{方}{かた}は\R{同}{おな}じです。\R{分母}{ぶん|ぼ}が13なら、あまりは0〜12の13\R{通}{とお}り。どんな\R{分数}{ぶん|すう}でも、かならず\R{終}{お}わるか、くり\R{返}{かえ}すかのどちらかになります。

\noindent\textbf{すぐにはくり\R{返}{かえ}さないこともある}\quad $\dfrac{1}{6}=0.1666\cdots$ は、\R{小数第}{しょう|すう|だい}1\R{位}{い}の1のあとから6がくり\R{返}{かえ}します（$10\div6=1$ あまり4、$40\div6=6$ あまり4、……と、あまり4がつづく）。はじめの\R{何}{なん}けたかがちがっても、どこかから\R{先}{さき}がくり\R{返}{かえ}せば、くり\R{返}{かえ}す\R{小数}{しょう|すう}のなかまです。

\begin{taisetsu}
\R{分数}{ぶん|すう}を\R{小数}{しょう|すう}にすると、\R{終}{お}わる\R{小数}{しょう|すう}か、くり\R{返}{かえ}す\R{小数}{しょう|すう}になる。
\end{taisetsu}

\bkiijhead{\R{反対}{はん|たい}に、\R{小数}{しょう|すう}を\R{分数}{ぶん|すう}にもどす}
\R{終}{お}わる\R{小数}{しょう|すう}は、$0.25=\dfrac{25}{100}=\dfrac{1}{4}$ のように、10・100・1000……を\R{分母}{ぶん|ぼ}にして\R{分数}{ぶん|すう}にもどせます。0.25は0.01が25こ\R{分}{ぶん}、つまり $\dfrac{1}{100}$ が25こ\R{分}{ぶん}だからです。

くり\R{返}{かえ}す\R{小数}{しょう|すう}はどうでしょう。まず $\dfrac{1}{9}$ を\R{小数}{しょう|すう}にします。$10\div9=1$ あまり1 で、あまりがいつも1なので、$\dfrac{1}{9}=0.111\cdots$ と1がずっとつづきます。\R{同}{おな}じように、$\dfrac{1}{99}$ は $100\div99=1$ あまり1 がつづくので、$\dfrac{1}{99}=0.010101\cdots$ です。
\begin{center}
\begin{tabular}{l@{\qquad}l}
$0.333\cdots$ は $0.111\cdots$ の3つ\R{分}{ぶん} & → $3\times\dfrac{1}{9}=\dfrac{3}{9}=\dfrac{1}{3}$ \\[8pt]
$0.121212\cdots$ は $0.010101\cdots$ の12こ\R{分}{ぶん} & → $12\times\dfrac{1}{99}=\dfrac{12}{99}=\dfrac{4}{33}$ \\
\end{tabular}
\end{center}
だから、\R{小数点}{しょう|すう|てん}のすぐ\R{後}{あと}からくり\R{返}{かえ}す\R{小数}{しょう|すう}は、くり\R{返}{かえ}す\R{部分}{ぶ|ぶん}を、くり\R{返}{かえ}すけたの\R{数}{かず}だけ9をならべた\R{数}{かず}（1けたなら9、2けたなら99）でわった\R{分数}{ぶん|すう}にもどせます。\R{電卓}{でん|たく}で $12\div99$ を\R{計算}{けい|さん}して\R{確}{たし}かめよう。\R{整数}{せい|すう}も $3=\dfrac{3}{1}$ と\R{書}{か}けます。このように、\R{分子}{ぶん|し}も\R{分母}{ぶん|ぼ}も\R{整数}{せい|すう}の\R{分数}{ぶん|すう}で\R{表}{あらわ}せる\R{数}{かず}を\textbf{\R{有理数}{ゆう|り|すう}}といいます。\bkidx{ゆうりすう}{\R{有理数}{ゆう|り|すう}}

\bksubsection{2}{$\sqrt{2}$ は\R{分数}{ぶん|すう}で\R{書}{か}けない}

\leavevmode\label{bk:2-9:muri}$\sqrt{2}$ に\R{近}{ちか}い\R{分数}{ぶん|すう}を2\R{乗}{じょう}して、ぴったり2になるものをさがします。たとえば $\dfrac{7}{5}$ を2\R{乗}{じょう}すると、\R{分数}{ぶん|すう}のかけ\R{算}{ざん}（2-4）で
\[ \frac{7}{5}\times\frac{7}{5}=\frac{7\times7}{5\times5}=\frac{49}{25} \]
2を\R{分母}{ぶん|ぼ}25の\R{分数}{ぶん|すう}にすると $\dfrac{50}{25}$ です。49と50はおしいけれど、ちがいますね。2\R{乗}{じょう}がぴったり2になるには、\textbf{\R{分子}{ぶん|し}×\R{分子}{ぶん|し}が、\R{分母}{ぶん|ぼ}×\R{分母}{ぶん|ぼ}の2\R{倍}{ばい}とぴったり\R{同}{おな}じ}でなければなりません。もっと\R{大}{おお}きな\R{分数}{ぶん|すう}でもためしてみましょう。

\begin{bkcenter}
\begin{bktable}{colspec={Q[c,m,wd=20mm] Q[c,m,wd=26mm] Q[c,m,wd=34mm] Q[c,m,wd=22mm]}, rowsep=1.2pt, colsep=4pt, cells={font=\small}, row{1}={font=\sffamily\bfseries\small}}
\R{分数}{ぶん|すう} & \R{分子}{ぶん|し}×\R{分子}{ぶん|し} & \R{分母}{ぶん|ぼ}×\R{分母}{ぶん|ぼ}×2 & 2\R{乗}{じょう}すると \\
$\dfrac{7}{5}$ & 49 & 50 & 1.96 \\
$\dfrac{17}{12}$ & 289 & 288 & 2.0069… \\
$\dfrac{41}{29}$ & 1681 & 1682 & 1.9988… \\
$\dfrac{99}{70}$ & 9801 & 9800 & 2.0002… \\
$\dfrac{577}{408}$ & 332929 & 332928 & 2.000006… \\
\end{bktable}
\bktabcaption{$\sqrt{2}$ に\R{近}{ちか}い\R{分数}{ぶん|すう}。いつも1だけずれて、ぴったり2にならない。}
\end{bkcenter}

じつは、どんな\R{分数}{ぶん|すう}を2\R{乗}{じょう}しても、ぴったり2にはならないことが\R{証明}{しょう|めい}できます（もっと\R{知}{し}りたい）。\R{分子}{ぶん|し}も\R{分母}{ぶん|ぼ}も\R{整数}{せい|すう}の\R{分数}{ぶん|すう}で\R{表}{あらわ}せない\R{数}{かず}を\textbf{\R{無理数}{む|り|すう}}といいます。\bkidx{むりすう}{\R{無理数}{む|り|すう}}$\sqrt{2}$ は\R{無理数}{む|り|すう}なので、\R{小数}{しょう|すう}にすると、\R{終}{お}わりもくり\R{返}{かえ}しもしません。これが\bkhatenaref の\R{答}{こた}えです。$\sqrt{\phantom{2}}$ の\R{中}{なか}が1、4、9、16、……（\R{整数}{せい|すう}の2\R{乗}{じょう}）でない\R{整数}{せい|すう}のときも、\R{同}{おな}じように\R{分数}{ぶん|すう}で\R{書}{か}けません。

\begin{motto}[3]
\textbf{$\sqrt{2}$ が\R{分数}{ぶん|すう}で\R{書}{か}けないわけ}\par
2でわりきれる\R{整数}{せい|すう}（2、4、6、……）を\textbf{\R{偶数}{ぐう|すう}}、わりきれない\R{整数}{せい|すう}（1、3、5、……）を\textbf{\R{奇数}{き|すう}}といいます。\R{奇数}{き|すう}×\R{奇数}{き|すう}は\R{奇数}{き|すう}（$3\times3=9$、$5\times7=35$）、\R{偶数}{ぐう|すう}にどんな\R{整数}{せい|すう}をかけても\R{偶数}{ぐう|すう}です。

もし $\sqrt{2}$ が、これ\R{以上}{い|じょう}\R{約分}{やく|ぶん}できない\R{分数}{ぶん|すう}で\R{書}{か}けたとすると、
\begin{enumerate}[label={\arabic*}.,leftmargin=1.6\zw, itemsep=0pt, topsep=2pt]
\item \R{分子}{ぶん|し}×\R{分子}{ぶん|し}＝2×\R{分母}{ぶん|ぼ}×\R{分母}{ぶん|ぼ}。\R{右}{みぎ}は\R{偶数}{ぐう|すう}なので、\R{分子}{ぶん|し}×\R{分子}{ぶん|し}も\R{偶数}{ぐう|すう}。\R{奇数}{き|すう}×\R{奇数}{き|すう}は\R{奇数}{き|すう}なので、\R{分子}{ぶん|し}は\R{偶数}{ぐう|すう}。
\item \R{分子}{ぶん|し}＝2×□と\R{書}{か}くと、4×□×□＝2×\R{分母}{ぶん|ぼ}×\R{分母}{ぶん|ぼ}。\R{両方}{りょう|ほう}を2でわって、\R{分母}{ぶん|ぼ}×\R{分母}{ぶん|ぼ}＝2×□×□。1.と\R{同}{おな}じように、\R{分母}{ぶん|ぼ}も\R{偶数}{ぐう|すう}。
\item \R{分子}{ぶん|し}も\R{分母}{ぶん|ぼ}も\R{偶数}{ぐう|すう}なら、2で\R{約分}{やく|ぶん}できます。「これ\R{以上}{い|じょう}\R{約分}{やく|ぶん}できない」に\R{合}{あ}いません。だから、「\R{分数}{ぶん|すう}で\R{書}{か}けた」がまちがいです。
\end{enumerate}
\end{motto}

\bksubsection{3}{\R{円周率}{えん|しゅう|りつ} $\pi$}

\leavevmode\label{bk:2-9:pi}\bkyattemiyouref では、\R{円}{えん}の\R{大}{おお}きさがちがっても、\R{円周}{えん|しゅう}÷\R{直径}{ちょっ|けい}はいつも3.14くらいになりました。この\R{数}{かず}を\textbf{\R{円周率}{えん|しゅう|りつ}}といい、ギリシャ\R{文字}{も|じ}の $\pi$（パイ）で\R{表}{あらわ}します。\bkidx{えんしゅうりつ}{\R{円周率}{えん|しゅう|りつ}（$\pi$）}
\[ \pi=3.14159265358979\cdots,\qquad \text{\R{円周}{えん|しゅう}＝\R{直径}{ちょっ|けい}×}\pi \]

\noindent\begin{minipage}[c]{0.6\linewidth}
はかると、どうしても\R{少}{すこ}しずれます。\R{計算}{けい|さん}で $\pi$ をはさんでみましょう。\R{円}{えん}の\R{内側}{うち|がわ}の\R{正六角形}{せい|ろっ|かく|けい}は、6つの\R{正三角形}{せい|さん|かく|けい}に\R{分}{わ}けられ、1\R{辺}{ぺん}が\R{半径}{はん|けい}と\R{同}{おな}じなので、まわりの\R{長}{なが}さは\R{半径}{はん|けい}6つ\R{分}{ぶん}、つまり\R{直径}{ちょっ|けい}の3\R{倍}{ばい}。\R{外側}{そと|がわ}の\R{正方形}{せい|ほう|けい}のまわりは\R{直径}{ちょっ|けい}の4\R{倍}{ばい}。\R{円周}{えん|しゅう}はその\R{間}{あいだ}なので、$\pi$ は3より\R{大}{おお}きく4より\R{小}{ちい}さい。\R{外側}{そと|がわ}も\R{正六角形}{せい|ろっ|かく|けい}にすると3.464…\R{倍}{ばい}。\R{角}{かど}の\R{数}{かず}をふやすと、\R{多角形}{た|かく|けい}が\R{円}{えん}にぴったりよりそっていき、はさむはばがせまくなります。
\end{minipage}\hfill
\begin{minipage}[c]{0.37\linewidth}\centering
\begin{tikzpicture}[x=1mm, y=1mm]
  \draw[bkTheme!80!black, line width=0.8pt] (-15,-15) rectangle (15,15);
  \draw[bkBlue, line width=1.2pt] (0,0) circle (15);
  \draw[bkGreen!70!black, line width=0.8pt] (0:15) -- (60:15) -- (120:15) -- (180:15) -- (240:15) -- (300:15) -- cycle;
  \draw[bkGray, line width=0.5pt] (0,0) -- (0:15) (0,0) -- (60:15);
  \fill[bkInk] (0,0) circle (0.6);
  \node[font=\sffamily\footnotesize, text=bkGreen!60!black] at (6.5,-4) {\R{半径}{はん|けい}};
  \node[anchor=north, font=\sffamily\footnotesize, text=bkInk] at (0,-17) {\R{内側}{うち|がわ}に\R{正六角形}{せい|ろっ|かく|けい}、\R{外側}{そと|がわ}に\R{正方形}{せい|ほう|けい}};
\end{tikzpicture}
\end{minipage}

\begin{bkcenter}
\begin{bktable}{colspec={Q[l,m,wd=33mm] *{5}{Q[c,m,wd=16mm]}}, rowsep=1.5pt, colsep=3pt, cells={font=\small}, row{1}={font=\sffamily\bfseries\small}}
\R{正多角形}{せい|た|かく|けい}の\R{角}{かど}の\R{数}{かず} & 6 & 12 & 24 & 48 & 96 \\
$\pi$ より\R{小}{ちい}さい（\R{内側}{うち|がわ}） & 3 & 3.105 & 3.132 & 3.139 & 3.1410 \\
$\pi$ より\R{大}{おお}きい（\R{外側}{そと|がわ}） & 3.465 & 3.216 & 3.160 & 3.147 & 3.1428 \\
\end{bktable}
\bktabcaption{\R{正多角形}{せい|た|かく|けい}のまわり÷\R{直径}{ちょっ|けい}。\R{内側}{うち|がわ}は\R{小}{ちい}さめに、\R{外側}{そと|がわ}は\R{大}{おお}きめにけたを\R{切}{き}った。2200\R{年}{ねん}ほど\R{前}{まえ}、アルキメデスは\R{正}{せい}96\R{角形}{かく|けい}で3.14まで\R{求}{もと}めた。}
\end{bkcenter}

12\R{角形}{かく|けい}から\R{先}{さき}の\R{数}{かず}は、4-4で\R{学}{まな}ぶ\R{三平方}{さん|へい|ほう}の\R{定理}{てい|り}をくり\R{返}{かえ}し\R{使}{つか}うと\R{計算}{けい|さん}できます。

\bkiijhead{なぜ、どんな\R{円}{えん}でも\R{円周}{えん|しゅう}÷\R{直径}{ちょっ|けい}は\R{同}{おな}じなの？}
\R{円}{えん}は、どれも\R{同}{おな}じ\R{形}{かたち}で、\R{大}{おお}きさだけがちがいます。\R{直径}{ちょっ|けい}を2\R{倍}{ばい}にした\R{円}{えん}は、\R{小}{ちい}さい\R{円}{えん}をそのまま2\R{倍}{ばい}に\R{拡大}{かく|だい}したものなので、まわりの\R{長}{なが}さも2\R{倍}{ばい}になります。\R{直径}{ちょっ|けい}も\R{円周}{えん|しゅう}も2\R{倍}{ばい}なので、わった\R{答}{こた}えは\R{変}{か}わりません（$\dfrac{2\times\text{\R{円周}{えん|しゅう}}}{2\times\text{\R{直径}{ちょっ|けい}}}$ は2で\R{約分}{やく|ぶん}できる）。\bkyattemiyouref で、かんでも\R{皿}{さら}でも3.14くらいになったのは、このためです。\R{円周}{えん|しゅう}÷\R{直径}{ちょっ|けい}が $\pi$ なので、かけ\R{算}{ざん}にもどすと、\R{円周}{えん|しゅう}＝\R{直径}{ちょっ|けい}×$\pi$ です（\R{速}{はや}さ＝\R{道}{みち}のり÷\R{時間}{じ|かん}から、\R{道}{みち}のり＝\R{速}{はや}さ×\R{時間}{じ|かん}にもどすのと\R{同}{おな}じ。2-3）。

\R{円周}{えん|しゅう}÷\R{直径}{ちょっ|けい}はいつも\R{同}{おな}じ\R{数}{かず}ですが、\R{円周}{えん|しゅう}と\R{直径}{ちょっ|けい}を\R{両方}{りょう|ほう}とも\R{整数}{せい|すう}にすることはできません。$\pi$ も\R{無理数}{む|り|すう}であることが\R{証明}{しょう|めい}されています（\R{大学}{だい|がく}で\R{学}{まな}ぶ\R{方法}{ほう|ほう}を\R{使}{つか}います）。3.14は、$\pi$ に\R{近}{ちか}い\R{数}{かず}です。

\bksubsection{4}{\R{数直線}{すう|ちょく|せん}は、すきまなく\R{数}{かず}でうまっている}

\leavevmode\label{bk:2-9:suchoku}\R{数直線}{すう|ちょく|せん}の\R{上}{うえ}に、$\sqrt{2}$ の\R{点}{てん}をコンパスでとってみます。1\R{辺}{ぺん}1の\R{正方形}{せい|ほう|けい}の\R{対角線}{たい|かく|せん}は $\sqrt{2}$ です（2-8の $\sqrt{\phantom{2}}$ のカードのとなりの\R{図}{ず}）。0を\R{中心}{ちゅう|しん}に、\R{対角線}{たい|かく|せん}の\R{長}{なが}さで\R{円}{えん}をかくと、\R{数直線}{すう|ちょく|せん}と $\sqrt{2}$ の\R{所}{ところ}で\R{交}{まじ}わります。

\begin{bkfigure}
\begin{tikzpicture}[x=1mm, y=1mm]
  \filldraw[fill=bkBlue!15, draw=bkBlue!70!black, line width=0.6pt] (0,0) rectangle (28,28);
  \draw[bkTheme!80!black, line width=0.9pt] (0,0) -- (28,28);
  \draw[bkTheme!80!black, line width=0.7pt, dashed] (45:39.6) arc (45:0:39.6);
  \draw[bk axis] (-6,0) -- (122,0);
  \foreach \x/\t in {0/0,28/1,56/2,84/3,112/4} {\draw[bkInk, line width=0.6pt] (\x,-1.4) -- (\x,1.4); \node[anchor=north, bk tick, inner sep=0pt] at (\x,-2.4) {\t};}
  \foreach \x/\t/\c in {14/{\tfrac{1}{2}}/bkGreen, 37.33/{\tfrac{4}{3}}/bkGreen, 39.6/{\sqrt{2}}/bkTheme, 76.11/{e}/bkPurple, 87.96/{\pi}/bkBlue} {
    \fill[\c] (\x,0) circle (1.1);}
  \node[anchor=south, font=\small, text=bkGreen!60!black, inner sep=0pt] at (14,2.4) {$\frac{1}{2}$};
  \node[anchor=north, font=\small, text=bkGreen!60!black, inner sep=0pt] at (35.6,-2.4) {$\frac{4}{3}$};
  \node[anchor=north west, font=\small, text=bkTheme!80!black, inner sep=0pt] at (40.6,-2.4) {$\sqrt{2}$};
  \node[anchor=south, font=\small, text=bkPurple, inner sep=0pt] at (76.11,2.4) {$e$};
  \node[anchor=south, font=\small, text=bkBlue, inner sep=0pt] at (87.96,2.4) {$\pi$};
  \node[font=\sffamily\footnotesize, text=bkTheme!80!black, align=center, inner sep=0pt] at (7.5,21) {\R{対角線}{たい|かく|せん}\\$=\sqrt{2}$};
\end{tikzpicture}
\bkfigcaption{\R{緑}{みどり}は\R{有理数}{ゆう|り|すう}、ほかは\R{無理数}{む|り|すう}。どちらも\R{数直線}{すう|ちょく|せん}の\R{上}{うえ}にある。}
\end{bkfigure}

\R{分数}{ぶん|すう}の\R{点}{てん}だけでは、$\sqrt{2}$ や $\pi$ の\R{所}{ところ}が「\R{穴}{あな}」になってしまいます。\R{有理数}{ゆう|り|すう}と\R{無理数}{む|り|すう}をあわせると、\R{数直線}{すう|ちょく|せん}はすきまなくうまります。5-14で\R{学}{まな}ぶ $e=2.71828\cdots$ も、\R{分数}{ぶん|すう}で\R{表}{あらわ}せない\R{数}{かず}です。\R{分数}{ぶん|すう}で\R{書}{か}けない\R{数}{かず}は、じつはたくさんあるのです。

\begin{taisetsu}
\R{分子}{ぶん|し}も\R{分母}{ぶん|ぼ}も\R{整数}{せい|すう}の\R{分数}{ぶん|すう}で\R{表}{あらわ}せる\R{数}{かず}が\R{有理数}{ゆう|り|すう}、\R{表}{あらわ}せない\R{数}{かず}が\R{無理数}{む|り|すう}。$\sqrt{2}$、$\pi$ は\R{無理数}{む|り|すう}。\\
\R{円周}{えん|しゅう}＝\R{直径}{ちょっ|けい}×$\pi$\qquad（$\pi=3.14159\cdots$、\R{計算}{けい|さん}では3.14を\R{使}{つか}う）
\end{taisetsu}

\begin{kotoba}
\small
\bkword{\R{有理数}{ゆう|り|すう}}{\R{分子}{ぶん|し}も\R{分母}{ぶん|ぼ}も\R{整数}{せい|すう}の\R{分数}{ぶん|すう}で\R{表}{あらわ}せる\R{数}{かず}。\R{小数}{しょう|すう}にすると、\R{終}{お}わるか、くり\R{返}{かえ}す。}
\bkword{\R{無理数}{む|り|すう}}{\R{分子}{ぶん|し}も\R{分母}{ぶん|ぼ}も\R{整数}{せい|すう}の\R{分数}{ぶん|すう}で\R{表}{あらわ}せない\R{数}{かず}。\R{小数}{しょう|すう}にすると、\R{終}{お}わらず、くり\R{返}{かえ}しもしない。}
\bkword{\R{円周率}{えん|しゅう|りつ}}{\R{円周}{えん|しゅう}÷\R{直径}{ちょっ|けい}。どんな\R{円}{えん}でも\R{同}{おな}じ\R{数}{かず}で、$\pi$ と\R{書}{か}く。}
\end{kotoba}

\bkkigocards{\bkkigocard{$\pi$}{パイ}{\R{円周率}{えん|しゅう|りつ}。\R{円周}{えん|しゅう}÷\R{直径}{ちょっ|けい}。$3.14159\cdots$}{pi}{2-9}}{%
\begin{tikzpicture}[x=1mm, y=1mm]
  \draw[bkBlue, line width=1pt] (0,6) circle (6);
  \draw[bkInk, line width=0.5pt] (-6,6) -- (6,6);
  \node[anchor=south, font=\sffamily\footnotesize, inner sep=1pt] at (0,6.3) {1};
  \draw[bkInk!70, line width=0.8pt] (0,0) -- (40,0);
  \draw[bkBlue, line width=2pt] (0,0) -- (37.7,0);
  \foreach \k in {0,1,2,3} \draw[bkInk, line width=0.6pt] (\k*12,-1.4) -- (\k*12,1.4);
  \foreach \k in {1,2,3} \node[anchor=north, font=\sffamily\footnotesize, inner sep=0pt] at (\k*12,-2.2) {\k};
  \node[anchor=south, font=\sffamily\footnotesize, text=bkBlue!80!black, inner sep=0pt] at (37.7,2.4) {$\pi$};
  \node[anchor=north, font=\sffamily\footnotesize, align=center] at (20,-6.5) {\R{直径}{ちょっ|けい}1の\R{円}{えん}が1\R{回転}{かい|てん}\\すると、$\pi$ だけ\R{進}{すす}む};
\end{tikzpicture}}

%% ---------------- 例題 ----------------
\bksection{reidai}{\R{例題}{れい|だい}でたしかめよう}

\begin{reidai}[\R{分数}{ぶん|すう}を\R{小数}{しょう|すう}にする]
（1）$\dfrac{3}{8}$\quad（2）$\dfrac{5}{11}$ を\R{小数}{しょう|すう}に。（3）$\dfrac{5}{11}$ の\R{小数}{しょう|すう}\R{第}{だい}10\R{位}{い}の\R{数字}{すう|じ}は？\par
\kangaekata \R{分子}{ぶん|し}÷\R{分母}{ぶん|ぼ}を\R{計算}{けい|さん}します。くり\R{返}{かえ}すときは、くり\R{返}{かえ}す\R{部分}{ぶ|ぶん}を\R{見}{み}つけます。\par
\shiki （1）$3\div8=0.375$\quad（あまりが6、4、0となって\R{終}{お}わる）\par
\hspace*{2.6\zw}（2）$5\div11=0.454545\cdots$\quad（あまりが6、5、6、5、……）\par
\hspace*{2.6\zw}（3）「45」の2けたを1\R{組}{くみ}と\R{考}{かんが}える。$10\div2=5$ あまり0\par
\hspace*{2.6\zw}\phantom{（3）}ちょうど5\R{組}{くみ}で\R{終}{お}わるので、\R{組}{くみ}の\R{最後}{さい|ご}のけたの5\par
\kotae （1）0.375（\R{終}{お}わる）\quad（2）$0.4545\cdots$（くり\R{返}{かえ}す）\quad（3）5
\end{reidai}

\begin{jibun}
（1）$\dfrac{7}{20}$\quad（2）$\dfrac{2}{7}$ を\R{小数}{しょう|すう}に。（3）$\dfrac{3}{11}$ の\R{小数}{しょう|すう}\R{第}{だい}9\R{位}{い}の\R{数字}{すう|じ}は？
\bkjans{（1）\nobreak 0.35\quad（2）\nobreak $0.285714\,285714\cdots$\quad（3）\nobreak $0.2727\cdots$。「27」が2けたずつ。$9\div2=4$ あまり1 なので、\R{組}{くみ}の1けためで2}
\end{jibun}

\begin{reidai}[くり\R{返}{かえ}しが\R{長}{なが}いときの\R{第}{だい}□\R{位}{い}]
$\dfrac{2}{7}=0.285714\,285714\cdots$ の、\R{小数第}{しょう|すう|だい}10\R{位}{い}の\R{数字}{すう|じ}は\R{何}{なん}ですか。\par
\kangaekata 「285714」の6けたが1\R{組}{くみ}でくり\R{返}{かえ}します。10けたの\R{中}{なか}に6けたの\R{組}{くみ}がいくつあり、\R{何}{なん}けたあまるかを、わり\R{算}{ざん}で\R{求}{もと}めます。\par
\shiki ① $10\div6=1$ あまり4\par
\hspace*{2.6\zw}② 1\R{組}{くみ}（6けた）が\R{終}{お}わったあと、4けためまで\R{進}{すす}む\par
\hspace*{2.6\zw}③ 「2、8、5、7、1、4」の4けためは7\par
\kotae 7
\end{reidai}

\begin{jibun}
$\dfrac{3}{7}=0.428571\,428571\cdots$ の、\R{小数第}{しょう|すう|だい}8\R{位}{い}の\R{数字}{すう|じ}は\R{何}{なん}ですか。
\bkjans{$8\div6=1$ あまり2 なので、「428571」の2けためで2}
\end{jibun}

\begin{reidai}[\R{小数}{しょう|すう}を\R{分数}{ぶん|すう}にする]
（1）0.6\quad（2）0.125\quad（3）$0.3333\cdots$ を\R{分数}{ぶん|すう}にしましょう。\par
\kangaekata \R{終}{お}わる\R{小数}{しょう|すう}は10・100・1000を\R{分母}{ぶん|ぼ}に、すぐくり\R{返}{かえ}す\R{小数}{しょう|すう}は9や99を\R{分母}{ぶん|ぼ}にして\R{約分}{やく|ぶん}します。\par
\shiki （1）$0.6=\dfrac{6}{10}=\dfrac{3}{5}$\qquad（2）$0.125=\dfrac{125}{1000}=\dfrac{1}{8}$\par
\hspace*{2.6\zw}（3）$0.333\cdots$ は $\dfrac{1}{9}=0.111\cdots$ の3つ\R{分}{ぶん}なので、$\dfrac{3}{9}=\dfrac{1}{3}$\par
\kotae （1）$\dfrac{3}{5}$\quad（2）$\dfrac{1}{8}$\quad（3）$\dfrac{1}{3}$
\end{reidai}

\begin{jibun}
（1）0.4\quad（2）0.45\quad（3）$0.1212\cdots$ を\R{分数}{ぶん|すう}にしましょう。
\bkjans{（1）\nobreak $\dfrac{4}{10}=\dfrac{2}{5}$\quad（2）\nobreak $\dfrac{45}{100}=\dfrac{9}{20}$\quad（3）\nobreak $\dfrac{12}{99}=\dfrac{4}{33}$}
\end{jibun}

\begin{reidai}[\R{有理数}{ゆう|り|すう}と\R{無理数}{む|り|すう}に\R{分}{わ}ける]
\R{次}{つぎ}の\R{数}{かず}を、\R{有理数}{ゆう|り|すう}と\R{無理数}{む|り|すう}に\R{分}{わ}けましょう。\quad 0.5、\ $\dfrac{2}{3}$、\ $\sqrt{9}$、\ $\sqrt{5}$、\ $\pi$、\ $-4$、\ $0.444\cdots$\par
\kangaekata \R{分子}{ぶん|し}も\R{分母}{ぶん|ぼ}も\R{整数}{せい|すう}の\R{分数}{ぶん|すう}で\R{書}{か}けるかどうかを\R{考}{かんが}えます。$\sqrt{\phantom{2}}$ は、\R{中}{なか}の\R{数}{かず}を\R{計算}{けい|さん}してから\R{判断}{はん|だん}します。\par
\shiki $0.5=\dfrac{1}{2}$、\ $\sqrt{9}=3=\dfrac{3}{1}$、\ $-4=\dfrac{-4}{1}$、\ $0.444\cdots=\dfrac{4}{9}$\par
\hspace*{2.6\zw}5は\R{整数}{せい|すう}を2\R{乗}{じょう}した\R{数}{かず}ではないので、$\sqrt{5}$ は\R{分数}{ぶん|すう}で\R{書}{か}けない（$\sqrt{2}$ と\R{同}{おな}じ\R{理由}{り|ゆう}。\pageref{bk:2-9:muri}ページ）\par
\hspace*{2.6\zw}$\pi$ は\R{無理数}{む|り|すう}（\pageref{bk:2-9:pi}ページ）\par
\begin{center}
\begin{tikzpicture}[x=1mm, y=1mm]
  \draw[bkGreen!70!black, fill=bkGreen!10, line width=0.8pt, rounded corners=2mm] (0,0) rectangle (66,22);
  \node[anchor=north west, font=\sffamily\bfseries\footnotesize, text=bkGreen!50!black] at (2,21) {\R{有理数}{ゆう|り|すう}（\R{分数}{ぶん|すう}で\R{書}{か}ける）};
  \node[font=\small] at (11,8) {0.5}; \node[font=\small] at (24,8) {$\dfrac{2}{3}$}; \node[font=\small] at (37,8) {$\sqrt{9}$}; \node[font=\small] at (48,8) {$-4$}; \node[font=\small] at (59.5,8) {$0.44\cdots$};
  \draw[bkPurple!70!black, fill=bkPurple!8, line width=0.8pt, rounded corners=2mm] (72,0) rectangle (118,22);
  \node[anchor=north west, font=\sffamily\bfseries\footnotesize, text=bkPurple!70!black] at (74,21) {\R{無理数}{む|り|すう}（\R{書}{か}けない）};
  \node[font=\small] at (86,8) {$\sqrt{5}$}; \node[font=\small] at (103,8) {$\pi$};
\end{tikzpicture}
\end{center}
\kotae \R{有理数}{ゆう|り|すう}：0.5、$\dfrac{2}{3}$、$\sqrt{9}$、$-4$、$0.444\cdots$\qquad\R{無理数}{む|り|すう}：$\sqrt{5}$、$\pi$
\end{reidai}

\begin{jibun}
\R{次}{つぎ}の\R{数}{かず}を、\R{有理数}{ゆう|り|すう}と\R{無理数}{む|り|すう}に\R{分}{わ}けましょう。\quad 1.25、\ $\sqrt{16}$、\ $\sqrt{3}$、\ $\dfrac{1}{6}$、\ $0.777\cdots$
\bkjans{\R{有理数}{ゆう|り|すう}：1.25、$\sqrt{16}$（$=4$）、$\dfrac{1}{6}$、$0.777\cdots$\quad \R{無理数}{む|り|すう}：$\sqrt{3}$}
\end{jibun}

\begin{reidai}[\R{円周}{えん|しゅう}と\R{直径}{ちょっ|けい}]
$\pi$ を3.14として、（1）\R{直径}{ちょっ|けい}\qty{10}{cm}の\R{円}{えん}の\R{円周}{えん|しゅう}、（2）\R{半径}{はん|けい}\qty{4}{cm}の\R{円}{えん}の\R{円周}{えん|しゅう}、（3）\R{円周}{えん|しゅう}\qty{62.8}{cm}の\R{円}{えん}の\R{直径}{ちょっ|けい}を\R{求}{もと}めましょう。\par
\kangaekata \R{円周}{えん|しゅう}＝\R{直径}{ちょっ|けい}×$\pi$。\R{半径}{はん|けい}がわかっているときは、まず2\R{倍}{ばい}して\R{直径}{ちょっ|けい}にします。\R{直径}{ちょっ|けい}は、\R{円周}{えん|しゅう}を $\pi$ でわれば\R{求}{もと}められます（\R{直径}{ちょっ|けい}×3.14＝62.8 の\R{直径}{ちょっ|けい}は、62.8÷3.14）。\par
\shiki （1）$10\times3.14=31.4$\par
\hspace*{2.6\zw}（2）\R{直径}{ちょっ|けい}は $4\times2=8$（cm）。$8\times3.14=25.12$\par
\hspace*{2.6\zw}（3）$62.8\div3.14=20$\par
\begin{center}
\begin{tikzpicture}[x=1mm, y=1mm]
  \draw[bkBlue, line width=1pt] (0,0) circle (11);
  \fill[bkInk] (0,0) circle (0.6);
  \draw[bkOrange!80!black, line width=0.9pt] (0,0) -- (11,0);
  \node[anchor=south, font=\sffamily\footnotesize, text=bkOrange!80!black, inner sep=1pt] at (5.5,0.6) {\R{半径}{はん|けい}4};
  \draw[bkGreen!60!black, line width=0.9pt] (-11,-0.01) -- (0,-0.01);
  \node[anchor=north, font=\sffamily\footnotesize, text=bkGreen!50!black, inner sep=1pt] at (-5.5,-0.6) {4};
  \node[anchor=west, font=\sffamily\footnotesize, align=left, inner sep=0pt] at (15,3) {\R{直径}{ちょっ|けい}は\R{半径}{はん|けい}の2つ\R{分}{ぶん}\\$4\times2=8$（cm）};
  \node[anchor=west, font=\sffamily\footnotesize, text=bkBlue, inner sep=0pt] at (15,-7) {\R{円周}{えん|しゅう}＝$8\times3.14$};
\end{tikzpicture}
\end{center}
\kotae （1）\qty{31.4}{cm}\quad（2）\qty{25.12}{cm}\quad（3）\qty{20}{cm}
\end{reidai}

\begin{jibun}
$\pi$ を3.14として、（1）\R{半径}{はん|けい}\qty{33}{cm}のタイヤの\R{円周}{えん|しゅう}と、（2）まわりの\R{長}{なが}さが\qty{157}{cm}の\R{木}{き}の\R{直径}{ちょっ|けい}を\R{求}{もと}めましょう。
\bkjans{（1）\nobreak \R{直径}{ちょっ|けい}は\qty{66}{cm}。$66\times3.14=207.24$ で\qty{207.24}{cm}\quad（2）\nobreak $157\div3.14=50$ で\qty{50}{cm}}
\end{jibun}

\begin{column}{$\pi$ に\R{近}{ちか}い\R{分数}{ぶん|すう}}
$\pi$ は\R{分数}{ぶん|すう}で\R{書}{か}けませんが、\R{近}{ちか}い\R{分数}{ぶん|すう}ならあります。\R{電卓}{でん|たく}で\R{確}{たし}かめてみよう。
\begin{bkcenter}
\begin{bktable}{colspec={Q[c,m,wd=20mm] Q[l,m,wd=44mm] Q[l,m,wd=46mm]}, rowsep=1.5pt, colsep=4pt, row{1}={bg=bkNavy!8}}
\R{分数}{ぶん|すう} & \R{小数}{しょう|すう}にすると & $\pi=3.14159265\cdots$ とくらべると \\
$\dfrac{22}{7}$ & $3.142857\,142857\cdots$ & \R{小数第}{しょう|すう|だい}2\R{位}{い}まで\R{同}{おな}じ \\
$\dfrac{355}{113}$ & $3.14159292\cdots$ & \R{小数第}{しょう|すう|だい}6\R{位}{い}まで\R{同}{おな}じ \\
\end{bktable}
\end{bkcenter}
$\dfrac{355}{113}$ は、1500\R{年}{ねん}ほど\R{前}{まえ}に\R{中国}{ちゅう|ごく}の\R{祖沖之}{そ|ちゅう|し}が\R{見}{み}つけた\R{分数}{ぶん|すう}です。どちらも\R{分数}{ぶん|すう}なので\R{有理数}{ゆう|り|すう}。$\pi$ そのものではありません。
\end{column}

\begin{reidai}[タイヤが\R{何回}{なん|かい}も\R{回}{まわ}ると、\R{何}{なん}m\R{進}{すす}む？]
\R{直径}{ちょっ|けい}\qty{50}{cm}の\R{車輪}{しゃ|りん}が500\R{回転}{かい|てん}すると、\R{何}{なん}m\R{進}{すす}みますか（$\pi$ は3.14）。\par
\kangaekata 1\R{回転}{かい|てん}で\R{円周}{えん|しゅう}の\R{長}{なが}さだけ\R{進}{すす}みます（\pageref{bk:2-9:tire}ページ）。500\R{回転}{かい|てん}なら、その500\R{倍}{ばい}。\R{答}{こた}えはmで\R{聞}{き}かれているので、さいごにcmをmになおします。\qty{1}{m}は\qty{100}{cm}なので、cmの\R{数}{かず}を100でわります。\par
\shiki ① 1\R{回転}{かい|てん}：$50\times3.14=157$（cm）\par
\hspace*{2.6\zw}② 500\R{回転}{かい|てん}：$157\times500=78500$（cm）\par
\hspace*{2.6\zw}③ $78500\div100=785$（m）\par
\kotae \qty{785}{m}
\end{reidai}

\begin{jibun}
\R{直径}{ちょっ|けい}\qty{60}{cm}の\R{車輪}{しゃ|りん}が1000\R{回転}{かい|てん}すると、\R{何}{なん}m\R{進}{すす}みますか。それは\R{何}{なん}kmですか（$\pi$ は3.14）。
\bkjans{$60\times3.14=188.4$\,cm、$188.4\times1000=188400$\,cm。$188400\div100=1884$ で\qty{1884}{m}。\qty{1}{km}＝\qty{1000}{m}なので\qty{1.884}{km}}
\end{jibun}

%% ---------------- まちがえやすいところ・物理工学 ----------------
\bksection{tsumazuki}{まちがえやすいところをチェック}

\begin{tsumazuki}{$\pi$ を3.14ぴったりと\R{思}{おも}う}
\bkng{$\pi=3.14$}
\bkok{$\pi=3.14159\cdots$。3.14は、\R{計算}{けい|さん}に\R{使}{つか}う\R{近}{ちか}い\R{数}{かず}}
$\pi$ は\R{無理数}{む|り|すう}なので、\R{小数}{しょう|すう}はどこまでも\R{終}{お}わりません。
\end{tsumazuki}

\begin{tsumazuki}{$0.333\cdots$ を $\dfrac{3}{10}$ にしてしまう}
\bkng{$0.333\cdots=\dfrac{3}{10}$}
\bkok{$0.333\cdots=\dfrac{3}{9}=\dfrac{1}{3}$}
$\dfrac{3}{10}=0.3$ で\R{終}{お}わる\R{小数}{しょう|すう}です。3がずっとつづく\R{小数}{しょう|すう}は、0.3より\R{少}{すこ}し\R{大}{おお}きく、$\dfrac{1}{9}$ の3つ\R{分}{ぶん}です。
\end{tsumazuki}

\begin{tsumazuki}{あまりが0のときの\R{第}{だい}□\R{位}{い}をまちがえる}
$\dfrac{5}{11}=0.4545\cdots$ の\R{小数第}{しょう|すう|だい}10\R{位}{い}は……（$10\div2=5$ あまり0）
\bkng{あまり0なので、「45」の1けための4}
\bkok{あまり0は、5\R{組}{くみ}めがちょうど\R{終}{お}わったところ。\R{組}{くみ}の\R{最後}{さい|ご}のけたの5}
あまりが1なら\R{次}{つぎ}の\R{組}{くみ}の1けため、あまりが0なら\R{組}{くみ}の\R{最後}{さい|ご}のけたです。
\end{tsumazuki}

\begin{tsumazuki}{\R{半径}{はん|けい}に $\pi$ をかけてしまう}
\bkng{\R{円周}{えん|しゅう}＝\R{半径}{はん|けい}×$\pi$}
\bkok{\R{円周}{えん|しゅう}＝\R{直径}{ちょっ|けい}×$\pi$＝\R{半径}{はん|けい}×2×$\pi$}
\R{半径}{はん|けい}のときは、まず2\R{倍}{ばい}して\R{直径}{ちょっ|けい}にします。
\end{tsumazuki}

\begin{tsumazuki}{\R{円周}{えん|しゅう}から\R{直径}{ちょっ|けい}を\R{求}{もと}めるのに、$\pi$ をかけてしまう}
\R{円周}{えん|しゅう}が\qty{31.4}{cm}の\R{円}{えん}の\R{直径}{ちょっ|けい}は……
\bkng{$31.4\times3.14=98.596$（cm）}
\bkok{$31.4\div3.14=10$（cm）}
\R{直径}{ちょっ|けい}×3.14＝31.4 なので、\R{直径}{ちょっ|けい}は31.4を3.14でわって\R{求}{もと}めます。\R{直径}{ちょっ|けい}は\R{円周}{えん|しゅう}より\R{短}{みじか}いはずなので、98.596はおかしいと\R{気}{き}づけます。
\end{tsumazuki}

\begin{tsumazuki}{$\sqrt{\phantom{2}}$ がついていれば\R{無理数}{む|り|すう}だと\R{思}{おも}う}
\bkng{$\sqrt{4}$ は\R{無理数}{む|り|すう}}
\bkok{$\sqrt{4}=2$ なので\R{有理数}{ゆう|り|すう}}
$\sqrt{\phantom{2}}$ の\R{中}{なか}が4、9、16のような\R{数}{かず}なら、\R{整数}{せい|すう}になります。
\end{tsumazuki}

\begin{butsuri}
\par\addvspace{3pt}\textbf{タイヤ1\R{回転}{かい|てん}で\R{進}{すす}む\R{距離}{きょ|り}}\label{bk:2-9:tire}\par\noindent タイヤが1\R{回転}{かい|てん}すると、\R{円周}{えん|しゅう}の\R{長}{なが}さだけ\R{進}{すす}みます。\R{直径}{ちょっ|けい}\qty{66}{cm}なら $66\times3.14=207.24$ で\qty{207.24}{cm}、100でわって\R{約}{やく}\qty{2.07}{m}です。\R{自転車}{じ|てん|しゃ}のスピードメーターは、\R{車輪}{しゃ|りん}の\R{磁石}{じ|しゃく}がセンサーの\R{前}{まえ}を\R{通}{とお}る\R{回数}{かい|すう}を\R{数}{かぞ}え、\R{回数}{かい|すう}×\R{円周}{えん|しゅう}で\R{進}{すす}んだ\R{距離}{きょ|り}を\R{計算}{けい|さん}しています。
\begin{bkfigure}
\begin{tikzpicture}[x=1mm, y=1mm]
  \draw[bkInk!70, line width=0.8pt] (-4,0) -- (82,0);
  \begin{scope}[opacity=0.35]\bkiijwheel{(66.5,10)}{10}\end{scope}
  \bkiijwheel{(10,10)}{10}
  \fill[bkRed] (10,10) ++(50:7.6) circle (1.2);
  \fill[bkInk!80] (19.5,17.5) rectangle (22.5,22.5);
  \draw[bkInk!60, line width=0.4pt] (22.5,20) -- (25,22.5);
  \node[anchor=west, font=\sffamily\footnotesize] at (25,23.5) {センサー};
  \draw[bkRed!80!black, line width=0.4pt] (14.2,17) -- (8,25.5);
  \node[anchor=south east, font=\sffamily\footnotesize, text=bkRed!80!black, inner sep=1pt] at (8,25.5) {\R{磁石}{じ|しゃく}};
  \fill[bkRed] (10,0) circle (1.1);
  \fill[bkRed] (66.5,0) circle (1.1);
  \draw[bkBlue, line width=0.7pt, <->] (10,-3.5) -- (66.5,-3.5);
  \node[anchor=north, font=\sffamily\footnotesize] at (38.25,-5) {1\R{回転}{かい|てん}で\R{円周}{えん|しゅう}（\R{直径}{ちょっ|けい}×$\pi$）だけ\R{進}{すす}む};
\end{tikzpicture}
\end{bkfigure}


\par\addvspace{3pt}\textbf{ロボットの\R{車輪}{しゃ|りん}}\par\noindent \R{車輪}{しゃ|りん}で\R{走}{はし}るロボットは、モーターが\R{何回転}{なん|かい|てん}したかを\R{数}{かぞ}えるセンサーで、\R{走}{はし}った\R{距離}{きょ|り}を\R{計算}{けい|さん}します。\R{直径}{ちょっ|けい}\qty{5.6}{cm}の\R{車輪}{しゃ|りん}なら、1\R{回転}{かい|てん}で $5.6\times3.14=17.584$ で\R{約}{やく}\qty{17.6}{cm}。\qty{1}{m}（\qty{100}{cm}）\R{進}{すす}ませたいときは、$100\div17.584=5.68\cdots$ で、\R{約}{やく}5.7\R{回転}{かい|てん}させるようにプログラムします。

\par\addvspace{3pt}\textbf{コンピュータの $\pi$}\par\noindent \R{無理数}{む|り|すう}の\R{小数}{しょう|すう}は\R{終}{お}わらないので、コンピュータは\R{有限}{ゆう|げん}のけたで\R{近}{ちか}い\R{値}{あたい}を\R{使}{つか}います。ふつうは $\pi=3.141592653589793$（16けたほど）。けたを\R{切}{き}るときに\R{生}{う}まれる\R{小}{ちい}さなずれは、2-19で\R{学}{まな}びます。
\end{butsuri}

\clearpage
%% ---------------- 章末問題 ----------------
\bkendproblems
\bkproblemgroup{\R{基本}{き|ほん}}{2}

\bkproblem \R{次}{つぎ}の\R{分数}{ぶん|すう}を\R{小数}{しょう|すう}にしましょう。\R{終}{お}わるか、くり\R{返}{かえ}すかも\R{答}{こた}えましょう。\par
（1）$\dfrac{5}{8}$\quad（2）$\dfrac{4}{9}$\quad（3）$\dfrac{3}{11}$
\bkpans{（1）\nobreak 0.625（\R{終}{お}わる）\quad（2）\nobreak $0.444\cdots$（くり\R{返}{かえ}す）\quad（3）\nobreak $0.2727\cdots$（くり\R{返}{かえ}す）}

\bkproblem \R{次}{つぎ}の\R{小数}{しょう|すう}を\R{分数}{ぶん|すう}にしましょう。\quad（1）0.8\quad（2）0.36\quad（3）$0.777\cdots$
\bkpans{（1）\nobreak $\dfrac{4}{5}$\quad（2）\nobreak $\dfrac{9}{25}$\quad（3）\nobreak $\dfrac{7}{9}$}

\bkproblem \R{次}{つぎ}の\R{数}{かず}を、\R{有理数}{ゆう|り|すう}と\R{無理数}{む|り|すう}に\R{分}{わ}けましょう。\quad $\sqrt{2}$、\ $\dfrac{3}{7}$、\ $\sqrt{25}$、\ 0.6、\ $\pi$、\ $\sqrt{10}$
\bkpans{\R{有理数}{ゆう|り|すう}：$\dfrac{3}{7}$、$\sqrt{25}$（$=5$）、0.6\quad \R{無理数}{む|り|すう}：$\sqrt{2}$、$\pi$、$\sqrt{10}$}

\bkproblem $\pi$ を3.14として、\R{円周}{えん|しゅう}を\R{求}{もと}めましょう。\par
（1）\R{直径}{ちょっ|けい}\qty{8}{cm}の\R{円}{えん}\quad（2）\R{半径}{はん|けい}\qty{5}{cm}の\R{円}{えん}
\bkpans{（1）\nobreak \qty{25.12}{cm}\quad（2）\nobreak \R{直径}{ちょっ|けい}\qty{10}{cm}なので\qty{31.4}{cm}}

\bkproblem \R{円周}{えん|しゅう}が\qty{94.2}{cm}の\R{円}{えん}の\R{直径}{ちょっ|けい}を\R{求}{もと}めましょう（$\pi$ は3.14）。
\bkpans{$94.2\div3.14=30$ で\qty{30}{cm}}

\bkproblem \R{次}{つぎ}の\R{数}{かず}のうち、\R{有理数}{ゆう|り|すう}をえらんで\R{整数}{せい|すう}で\R{表}{あらわ}しましょう。\par
$\sqrt{36}$、\ $\sqrt{40}$、\ $\sqrt{49}$、\ $\sqrt{50}$
\bkpans{$\sqrt{36}=6$、$\sqrt{49}=7$（$\sqrt{40}$ と $\sqrt{50}$ は\R{無理数}{む|り|すう}）}

\bkproblemgroup{チャレンジ}{2}

\bkproblem $\dfrac{1}{7}=0.142857\,142857\cdots$ の、\R{小数}{しょう|すう}\R{第}{だい}20\R{位}{い}の\R{数字}{すう|じ}は\R{何}{なん}ですか。
\bkpans{6けたずつくり\R{返}{かえ}す。$20\div6=3$ あまり2 なので、くり\R{返}{かえ}しの2つ\R{目}{め}の4}

\bkproblem \R{直径}{ちょっ|けい}\qty{66}{cm}のタイヤが1000\R{回転}{かい|てん}すると、\R{約}{やく}\R{何}{なん}km\R{進}{すす}みますか（$\pi$ は3.14）。
\bkpans{$66\times3.14=207.24$\,cm。$207.24\times1000=207240$\,cm。$207240\div100=2072.4$\,m、$2072.4\div1000=2.0724$ で、\R{約}{やく}\qty{2.07}{km}}

\bkproblem $\dfrac{4}{7}=0.571428\,571428\cdots$ の、\R{小数第}{しょう|すう|だい}18\R{位}{い}の\R{数字}{すう|じ}は\R{何}{なん}ですか。
\bkpans{$18\div6=3$ あまり0 なので、3\R{組}{くみ}めがちょうど\R{終}{お}わる。「571428」の\R{最後}{さい|ご}の8}

\bkproblem $\pi$ の\R{代}{か}わりに $\dfrac{22}{7}$ を\R{使}{つか}うことがあります。$\dfrac{22}{7}$ を\R{小数}{しょう|すう}にして $\pi=3.14159\cdots$ とくらべましょう。$\dfrac{22}{7}$ は\R{有理数}{ゆう|り|すう}ですか、\R{無理数}{む|り|すう}ですか。
\bkpans{$\dfrac{22}{7}=3.142857\,142857\cdots$ で、$\pi$ より\R{少}{すこ}し\R{大}{おお}きい。\R{分子}{ぶん|し}も\R{分母}{ぶん|ぼ}も\R{整数}{せい|すう}の\R{分数}{ぶん|すう}なので\R{有理数}{ゆう|り|すう}}

\begin{column}{$\pi$ は、どこまで\R{求}{もと}められた？}
\begin{bkcenter}
\begin{bktable}{colspec={Q[l,m,wd=27mm] Q[l,m,wd=47mm] Q[l,m,wd=36mm]}, rowsep=1.5pt, colsep=4pt, row{1}={bg=bkNavy!8}}
いつ & だれが & どこまで \\
\R{約}{やく}2200\R{年}{ねん}\R{前}{まえ} & アルキメデス（ギリシャ） & 3.14（\R{正}{せい}96\R{角形}{かく|けい}で） \\
\R{約}{やく}340\R{年}{ねん}\R{前}{まえ} & \R{関孝和}{せき|たか|かず}（\R{江戸}{え|ど}\R{時代}{じ|だい}） & \R{小数}{しょう|すう}\R{第}{だい}10\R{位}{い}まで \\
2022\R{年}{ねん} & コンピュータ & 100\R{兆}{ちょう}けた（その\R{後}{ご}ものびている） \\
\end{bktable}
\end{bkcenter}
どこまで\R{計算}{けい|さん}しても、くり\R{返}{かえ}しは\R{出}{で}てきません。$\pi$ が\R{無理数}{む|り|すう}だからです。
\end{column}

\begin{furikaeri}
\begin{checklist}
\item \R{分数}{ぶん|すう}を\R{小数}{しょう|すう}にすると、\R{終}{お}わるかくり\R{返}{かえ}すわけがわかる。\hfill{\footnotesize\sffamily\color{bkTheme}→\pageref{bk:2-9:kurikaesu}ページ}
\item \R{分数}{ぶん|すう}で\R{書}{か}けない\R{数}{かず}（\R{無理数}{む|り|すう}）があることがわかる。\hfill{\footnotesize\sffamily\color{bkTheme}→\pageref{bk:2-9:muri}ページ}
\item \R{円周}{えん|しゅう}＝\R{直径}{ちょっ|けい}×$\pi$ で、\R{円周}{えん|しゅう}や\R{直径}{ちょっ|けい}を\R{求}{もと}められる。\hfill{\footnotesize\sffamily\color{bkTheme}→\pageref{bk:2-9:pi}ページ}
\end{checklist}
\end{furikaeri}

\bkfukidashi{delta}{\R{数直線}{すう|ちょく|せん}には、\R{分数}{ぶん|すう}で\R{書}{か}けない\R{数}{かず}もぎっしりならんでいるんだ。\R{次}{つぎ}の2-10では、はかった\R{量}{りょう}の「\R{単位}{たん|い}」の\R{世界}{せ|かい}ルールを\R{学}{まな}ぶよ。}
